\begin{statement}{C.2.RAT.6}{Known, with exact derivation}{Cubic suspensions of plane curves with delta at most five}
Let $r(x,y)=0$ be a reduced plane-curve germ with
$\delta(r)\le5$, and let $X_r$ be the isolated surface
singularity $z^3=r(x,y)$. Then $p_g(X_r)\le1$ and the
first rational Betti number of its link is at most two.
The latter is two precisely for plane type $A_5$ or $D_4$.
For a collection of these plane germs with total delta
invariant at most five,
\[
             \sum p_g(X_r)\le1,
           \qquad \sum b_1(\operatorname{link}X_r)\le2.
\tag{C.RD.5}
\]
\end{statement}
\begin{proof}
We first check that the finite normal-form list is exhaustive.
The delta formula under a point blowup shows that a plane
germ of multiplicity at least four has delta at least six.
For an irreducible branch of multiplicity two,
Weierstrass preparation, completion of the square and a
root of a unit give $y^2-x^{2k+1}$, of delta $k\le5$.
For an irreducible branch of multiplicity three, the
Puiseux conductor formula gives delta at least six
unless its first characteristic exponent is four or
five. A term with exponent divisible by three is removed
by subtracting a power of the first coordinate. The
remaining two possibilities have relatively prime
exponents $(3,4)$ or $(3,5)$, and no later characteristic
pair. Their analytic normal forms are $E_6$ and $E_8$.
This last analytic comparison is also the zero-parameter
case of the semiquasihomogeneous normal-form theorem:
the bases of the two gradient algebras are respectively
$x^i y^j$ with $i\le2,j\le1$, and with
$i\le3,j\le1$, after interchanging variables if necessary.
Every basis monomial has weight less than one, so there
is no higher-weight term left in the normal form; see
\cite[\S1, I.1.0, and \S2]{Arnold1976NormalForms}.

For a reducible germ the exact normalization sequence
for a union gives
\[
 \delta\Bigl(\bigcup_i D_i\Bigr)
       =\sum_i\delta(D_i)+\sum_{i<j}I(D_i,D_j).
\tag{C.RD.6}
\]
Four branches already contribute at least six.
With three branches, a singular branch would contribute
at least $1+2+2+1=6$. Thus all three are smooth.
They are either pairwise transverse, or exactly one
pair is tangent with contact two or three: if all
three had a common tangent, the sum would again be
at least six. Straightening the transverse branch,
subtracting the mean of the other two graph functions
and taking a root of their unit difference gives
$D_4,D_6,D_8$ respectively.

With two smooth branches, subtracting their mean graph
and rescaling the transverse coordinate gives
$A_{2k-1}$ with contact $k\le5$. With two singular
branches the two delta contributions and their
intersection multiplicity give at least $1+1+4=6$.
The remaining possibility is one smooth branch and
one cusp. A multiplicity-three cusp would already
give $3+3=6$. A multiplicity-two cusp of delta $k$
has parameterization $(t^2,t^{2k+1})$.
Its intersection with a smooth branch is either two,
or, for the common tangent, one of
$4,6,\ldots,2k,2k+1$.
The inequality $k+I\le5$ leaves a transverse smooth
branch with $k=1,2,3$, giving $D_5,D_7,D_9$, or the
ordinary cusp with its tangent smooth branch, giving
$E_7$. The actual analytic forms are again those
obtained by the same graph changes and square
completion, or the zero-parameter normal forms in the
cited theorem. The resulting complete list is
\[
              A_1,\ldots,A_{10},\qquad
              D_4,\ldots,D_9,\qquad E_6,E_7,E_8.
\tag{C.RD.7}
\]
An analytic equivalence of the plane equation up to
a unit gives an analytic equivalence of the cubic
suspensions: take a local cube root of that unit and
rescale $z$. Thus these are the precise surface germs
whose invariants must be computed.

The normal forms, with weights normalized so the
equation has degree one, are
\[
\begin{array}{c|c|c|c}
 \text{plane type}&r(x,y)&(w_x,w_y)&\delta(r)\\\hline
 A_n&x^2+y^{n+1}&(1/2,1/(n+1))&\lfloor(n+1)/2\rfloor\\
 D_n&x^2y+y^{n-1}&((n-2)/(2(n-1)),1/(n-1))
                         &\lfloor(n+2)/2\rfloor\\
 E_6&x^3+y^4&(1/3,1/4)&3\\
 E_7&x^3+xy^3&(1/3,2/9)&4\\
 E_8&x^3+y^5&(1/3,1/5)&4.
\end{array}
\tag{C.RD.8}
\]
The third weight of the suspension is $1/3$.
The monomial-spectrum formula and its geometric-genus
interval, or the direct adjunction-ideal calculation
\cite[Theorem 2.1.1 and Corollary 2.1.3]{MerleTeissier1977},
give
\[
 p_g(X_r)=\#\{(i,j,k)\in\mathbb Z_{\ge0}^3:
          (i+1)w_x+(j+1)w_y+(k+1)/3\le1\}.
\tag{C.RD.9}
\]
These particular forms are Newton nondegenerate;
this follows directly by checking that their face
polynomials have no common zero of their derivatives
in the relevant algebraic torus. In (C.RD.9) either
there is no monomial, for $A_1,\ldots,A_4$, or the
constant monomial alone satisfies the inequality.
In particular a nonzero geometric genus requires
delta at least three.

For clarity, the link calculation uses the actual
Milnor-fibre monodromy rather than a homology-sphere
assumption. For a monomial $x^i y^j$ in the plane
Jacobian algebra, its monodromy exponent modulo
integers is $(i+1)w_x+(j+1)w_y$. The two reduced
eigenvalues of $z^3$ are the two primitive cube
roots. The Thom--Sebastiani tensor comparison
therefore makes the eigenvalue-one dimension for
the suspension equal to the number of plane
exponents congruent to $1/3$ or $2/3$.
The isolated-surface Milnor fibre is simply
connected. Its pair sequence and Poincare--Lefschetz
duality identify the rational first homology of
its boundary with the radical of the middle
intersection pairing. The variation comparison
identifies that radical dimension with the
eigenvalue-one dimension. These weighted homogeneous
monodromies are semisimple; the monomial spectrum
and suspension formula are given in
\cite[pp. 165--167]{Steenbrink1989Spectrum}.

For $A_n$ the plane Jacobian basis is
$1,y,\ldots,y^{n-1}$; for $D_n$ it is
$1,y,\ldots,y^{n-2},x$; for $E_6,E_8$ it is the
rectangular basis from the preceding paragraph.
For $E_7$ it is $1,y,\ldots,y^4,x,xy$.
Direct substitution gives
\[
\begin{array}{c|c|c}
 \text{type}&p_g(z^3-r)&b_1(\operatorname{link}(z^3-r))\\\hline
 A_1,\ldots,A_4&0&0\\
 A_5&1&2\\
 A_6,\ldots,A_{10}&1&0\\
 D_4&1&2\\
 D_5,\ldots,D_9&1&0\\
 E_6,E_7,E_8&1&0.
\end{array}
\tag{C.RD.10}
\]
All weights, the individual newly required
$A_9,A_{10},D_8,D_9$ rows and every unordered
collection with total delta at most five are
checked with exact rational arithmetic in
\path{code/A_cubic_delta5.py} and its saved output.
Since each nonzero entry in either invariant
costs delta at least three, at most one such
entry occurs in a collection of total delta
at most five. This proves (C.RD.5).
\end{proof}


% C.2.RAT.7 from part_C2_rational_triple.tex

